\section{Anexo B: Transformadores Densos, Recurrentes y Aumentados con Memoria}
\label{sec:annex-transformer}

Este anexo exhaustivo sintetiza las arquitecturas modernas de transformers más allá de la línea base estándar de softmax. El campo ha evolucionado hacia seis paradigmas principales: atención densa global con variantes, compresión de cabezas clave-valor mediante atención por consultas agrupadas, compresión de estado recurrente con decaimiento, modelos de espacio de estado selectivos, integración numérica de profundidad continua, arquitecturas aumentadas con memoria y enfoques híbridos. Comprender esta taxonomía ilumina las compensaciones fundamentales entre expresividad, costo computacional, huella de memoria y simplicidad de despliegue.

\subsection{Atención Densa de Referencia: Estándar y Sigmoide}

\subsubsection{Atención Softmax Estándar}
La atención multi-cabeza estándar \citep{vaswani_attention_2017} calcula la similitud de producto punto escalado entre embeddings de tokens:

\begin{algorithmic}[1]
\State \textbf{Entrada:} Matriz de consulta $\mathbf{Q} \in \mathbb{R}^{n \times d}$, Matriz de clave $\mathbf{K} \in \mathbb{R}^{n \times d}$, Matriz de valor $\mathbf{V} \in \mathbb{R}^{n \times d}$
\State $\text{puntajes} \gets \frac{\mathbf{Q}\mathbf{K}^\top}{\sqrt{d}} \in \mathbb{R}^{n \times n}$ \Comment{calcular similitudes por pares}
\State $\text{pesos\_atencion} \gets \texttt{softmax}(\text{puntajes}, \text{eje}=1) \in \mathbb{R}^{n \times n}$ \Comment{normalizar sobre las claves}
\State \textbf{Salida:} $\mathbf{Y} = \text{pesos\_atencion} \cdot \mathbf{V} \in \mathbb{R}^{n \times d}$ \Comment{combinación ponderada de valores}
\end{algorithmic}

Para generación autorregresiva, el almacenamiento en caché KV guarda claves y valores pasados para evitar el recálculo $\mathcal{O}(n^2)$. Sin embargo, este caché de crecimiento lineal ocupa $\sim n \cdot d_{\text{oculto}}$ bytes, lo que puede consumir cientos de gigabytes para modelos de miles de millones de parámetros. La atención estándar logra una \textbf{expresividad perfecta} dentro de la ventana de contexto---cualquier token puede atender a cualquier otro token con pesos aprendidos---pero paga el precio del cómputo denso.

\begin{itemize}[leftmargin=*]
    \item \textbf{Complejidad de entrenamiento}: $\mathcal{O}(n^2 \cdot d)$ tiempo, $\mathcal{O}(n^2)$ espacio (materialización de la matriz de atención).
    \item \textbf{Complejidad de inferencia}: $\mathcal{O}(n)$ tiempo por token, $\mathcal{O}(n)$ espacio (caché KV).
    \item \textbf{Fortalezas}: Expresividad incomparable; recuperación perfecta del historial; altamente paralelizable.
    \item \textbf{Debilidades}: El cuello de botella cuadrático prohíbe contextos muy largos; el caché KV domina la memoria durante la generación.
\end{itemize}

\subsubsection{Atención Sigmoide}
La atención sigmoide reemplaza el softmax por filas con una activación sigmoide elemento a elemento \citep{ramapuram_theory_2024}:

\begin{algorithmic}[1]
\State \textbf{Entrada:} $\mathbf{Q}, \mathbf{K}, \mathbf{V}$ como arriba, sesgo aprendible $\mathbf{b} \in \mathbb{R}^{n \times n}$
\State $\text{logits} \gets \frac{\mathbf{Q}\mathbf{K}^\top}{\sqrt{d}} + \mathbf{b}$
\State $\text{pesos\_atencion} \gets \sigma(\text{logits})$ \Comment{sigmoide elemento a elemento, no softmax}
\State \textbf{Salida:} $\mathbf{Y} = \text{pesos\_atencion} \odot \mathbf{V}$
\end{algorithmic}

A diferencia de softmax, la sigmoide no impone una distribución de probabilidad (los pesos no necesitan sumar 1), lo que permite una independencia de tokens más fuerte. El análisis teórico mediante mezcla de expertos muestra que la sigmoide alcanza una complejidad muestral superior: convergencia $\mathcal{O}(n^{-0.51})$ para expertos ReLU frente a $\mathcal{O}(n^{-0.24})$ de softmax. Sin embargo, el entrenamiento empírico reveló inestabilidades de gradiente a escala. El remedio es la \textbf{norma híbrida}---agregar normalización después de la salida de atención---que estabiliza los gradientes sin sacrificar los beneficios teóricos del enrutamiento elemento a elemento.

\begin{itemize}[leftmargin=*]
    \item \textbf{Complejidad de entrenamiento}: $\mathcal{O}(n^2 \cdot d)$ (idéntico al estándar), pero las operaciones elemento a elemento permiten una aceleración de inferencia del 17\% mediante FlashSigmoid.
    \item \textbf{Complejidad de inferencia}: $\mathcal{O}(n)$ por token con caché KV (asintóticamente igual, pero con factores constantes más bajos).
    \item \textbf{Fortalezas}: Supera la competencia de suma cero; evita la sincronización por filas; implementación eficiente en hardware.
    \item \textbf{Debilidades}: Requiere estabilización cuidadosa (norma híbrida); inestabilidad de entrenamiento a grandes escalas sin pérdida auxiliar.
\end{itemize}

\subsubsection{Atención por Consultas Agrupadas (GQA)}
La atención por consultas agrupadas (GQA) \citep{ainslie_gqa_2023} tiende un puente entre la atención multi-cabeza (MHA) y la atención multi-consulta (MQA) al particionar las cabezas de consulta en grupos, cada uno compartiendo una sola cabeza clave-valor. Si $h$ es el número total de cabezas de atención, $h_k$ es el número de cabezas clave-valor, y $G = h / h_k$ es el tamaño del grupo, entonces cada grupo de $G$ cabezas de consulta comparte una cabeza KV:

\begin{algorithmic}[1]
\State \textbf{Entrada:} $\mathbf{Q} \in \mathbb{R}^{n \times h \cdot d_h}$, $\mathbf{K} \in \mathbb{R}^{n \times h_k \cdot d_h}$, $\mathbf{V} \in \mathbb{R}^{n \times h_k \cdot d_h}$
\State $\mathbf{K}_{\text{completo}} \gets \texttt{repeat\_interleave}(\mathbf{K}, G)$ \Comment{expandir de $h_k$ a $h$ cabezas}
\State $\mathbf{V}_{\text{completo}} \gets \texttt{repeat\_interleave}(\mathbf{V}, G)$
\State $\text{puntajes} \gets \frac{\mathbf{Q}\mathbf{K}_{\text{completo}}^\top}{\sqrt{d_h}} \in \mathbb{R}^{n \times n}$ \Comment{producto punto estándar tras expansión}
\State $\text{pesos\_atencion} \gets \texttt{softmax}(\text{puntajes}, \text{eje}=1)$
\State \textbf{Salida:} $\mathbf{Y} = \text{pesos\_atencion} \cdot \mathbf{V}_{\text{completo}} \in \mathbb{R}^{n \times h \cdot d_h}$
\end{algorithmic}

GQA interpola suavemente entre MHA ($h_k = h$, $G = 1$) y MQA ($h_k = 1$, $G = h$). La ventaja clave es un balance ajustable entre la huella del caché KV y la capacidad representacional. Cada cabeza KV adicional consume $2 \cdot d_h \cdot n$ bytes extra en el caché, por lo que reducir $h_k$ produce ahorros proporcionales de memoria.

\begin{itemize}[leftmargin=*]
    \item \textbf{Complejidad de entrenamiento}: $\mathcal{O}(n^2 \cdot h \cdot d_h)$ (idéntico a MHA tras la expansión de cabezas).
    \item \textbf{Complejidad de inferencia}: $\mathcal{O}(n)$ por token, espacio de caché KV $\mathcal{O}(h_k \cdot n \cdot d_h)$.
    \item \textbf{Fortalezas}: Balance flexible de cabezas; probado a escala (Llama 2 y 3 usan $h_k = 8$ cabezas KV con $h = 32$ cabezas de consulta); reduce el caché KV en factor $G \times$; soporta inferencia rápida con FlashAttention.
    \item \textbf{Debilidades}: Menos cabezas KV pueden reducir la capacidad representacional en modelos pequeños; requerir que $h_k$ divida a $h$ restringe las opciones arquitectónicas.
\end{itemize}

\subsection{Arquitecturas Recurrentes y Retentivas}

\subsubsection{Redes Retentivas (RetNet)}
RetNet \citep{sun_retentive_2023} unifica tres paradigmas de cómputo: entrenamiento paralelo, inferencia recurrente y despliegue por fragmentos. Su innovación central es el \textbf{mecanismo de retención}, que utiliza una matriz de decaimiento exponencial fija para modelar la importancia temporal:

\begin{algorithmic}[1]
\State \textbf{Forma paralela (entrenamiento):}
\State $\text{matriz\_decaimiento}[i,j] \gets \gamma^{i-j}$ para $i \geq j$, si no $0$ \Comment{decaimiento exponencial causal}
\State $\text{matriz\_decaimiento}[i,j] \gets 0$ para $i < j$ \Comment{enmascaramiento causal}
\State $\mathbf{Y}_{\text{paralelo}} \gets (\mathbf{Q}\mathbf{K}^\top \odot \text{matriz\_decaimiento}) \mathbf{V}$ \Comment{multiplicación elemento a elemento con decaimiento}
\end{algorithmic}

\begin{algorithmic}[1]
\State \textbf{Forma recurrente (inferencia):}
\For{$t = 1, \ldots, n$}
    \State $\mathbf{s}_t \gets \gamma \mathbf{s}_{t-1} + \mathbf{k}_t \mathbf{v}_t^\top$ \Comment{actualización de estado con decaimiento}
    \State $\mathbf{y}_t \gets \mathbf{q}_t \mathbf{s}_t$ \Comment{salida mediante interacción consulta-estado}
\EndFor
\end{algorithmic}

El escalar de decaimiento $\gamma \in (0,1)$ controla la ventana temporal. RetNet utiliza \textbf{retención multi-escala} con diferentes valores de $\gamma$ por cabeza (ej., $\gamma = 1 - 2^{-5}, 1 - 2^{-6}, \ldots$), permitiendo dependencias a corto y largo plazo simultáneamente. El modo recurrente por fragmentos divide las secuencias en fragmentos, procesa cada fragmento en paralelo y enhebra un estado recurrente entre fragmentos.

\begin{itemize}[leftmargin=*]
    \item \textbf{Complejidad de entrenamiento}: $\mathcal{O}(n^2 \cdot d)$ (paralela), o $\mathcal{O}(n \cdot c \cdot d)$ (recurrente por fragmentos con tamaño de fragmento $c$).
    \item \textbf{Complejidad de inferencia}: $\mathcal{O}(1)$ por token, $\mathcal{O}(d^2)$ espacio de estado (matriz fija).
    \item \textbf{Fortalezas}: Inferencia en tiempo constante; triple paradigma de cómputo; consolidación multi-escala.
    \item \textbf{Debilidades}: El decaimiento fijo impone un sesgo inductivo rígido; puede truncar patrones de largo alcance aprendidos.
\end{itemize}

\subsubsection{Mamba: Modelos de Espacio de Estado Selectivos}
Mamba \citep{gu_mamba_2023} enmarca la recurrencia como un sistema dinámico de tiempo continuo con parámetros dependientes de la entrada, logrando tanto entrenamiento lineal como inferencia en tiempo constante:

\begin{algorithmic}[1]
\State \textbf{Dinámica continua:} $h'(t) = \mathbf{A} h(t) + \mathbf{B} x(t)$, \quad $y(t) = \mathbf{C} h(t)$
\State \textbf{Discretización con tamaño de paso $\Delta_t$:}
\State $\bar{\mathbf{A}}_t \gets \exp(\Delta_t \mathbf{A})$
\State $\bar{\mathbf{B}}_t \gets (\Delta_t \mathbf{A})^{-1} (\exp(\Delta_t \mathbf{A}) - \mathbf{I}) \Delta_t \mathbf{B}_t$
\State \textbf{Actualización recurrente:}
\For{$t = 1, \ldots, n$}
    \State $\Delta_t \gets \text{softplus}(\text{Lineal}(x_t))$ \Comment{tamaño de paso dependiente de la entrada}
    \State $\mathbf{B}_t \gets \text{Lineal}(x_t)$, \quad $\mathbf{C}_t \gets \text{Lineal}(x_t)$ \Comment{matrices de proyección dependientes de la entrada}
    \State $h_t \gets \bar{\mathbf{A}}_t h_{t-1} + \bar{\mathbf{B}}_t x_t$ \Comment{transición de estado}
    \State $y_t \gets \mathbf{C}_t h_t$ \Comment{proyección de salida}
\EndFor
\end{algorithmic}

La innovación crítica es que $\mathbf{A}$ (la matriz del sistema) \textit{no} depende de la entrada, pero $\Delta_t$, $\mathbf{B}_t$ y $\mathbf{C}_t$ sí, haciendo que el sistema sea \textbf{variante en el tiempo}. Esta selectividad permite al modelo \textit{ignorar} información irrelevante estableciendo $\Delta_t$ cerca de cero, creando efectivamente una compuerta. Un algoritmo de escaneo paralelo consciente del hardware implementa la recurrencia eficientemente en GPUs fusionando el cómputo dentro de SRAM, evitando el costoso ancho de banda de HBM.

\begin{itemize}[leftmargin=*]
    \item \textbf{Complejidad de entrenamiento}: $\mathcal{O}(n \cdot d)$ mediante escaneo consciente del hardware (lineal en la longitud de la secuencia).
    \item \textbf{Complejidad de inferencia}: $\mathcal{O}(1)$ por token, $\mathcal{O}(d)$ estado (vector oculto).
    \item \textbf{Fortalezas}: Logra entrenamiento lineal e inferencia constante simultáneamente; práctico en secuencias largas (millones de tokens).
    \item \textbf{Debilidades}: La compresión del vector de estado puede debilitar la copia exacta y la recuperación asociativa densa en comparación con la atención completa.
\end{itemize}

\subsection{Transformers de Profundidad Continua: Integración EDO}

\subsubsection{Transformer EDO}
El Transformer EDO \citep{zhang_continuous_2021} interpreta la profundidad de la red como integración numérica de un sistema dinámico continuo, utilizando solucionadores Runge-Kutta de orden superior para reducir el error de truncamiento:

\begin{algorithmic}[1]
\State \textbf{Formulación continua:} $\frac{dh(t)}{dt} = f_\theta(h(t), t)$ donde $f_\theta$ es la subred del transformer
\State \textbf{Aproximación discreta Runge-Kutta-4:}
\State $k_1 \gets f_\theta(h_t, t)$
\State $k_2 \gets f_\theta(h_t + \tfrac{1}{2}k_1, t + \tfrac{1}{2}\Delta t)$
\State $k_3 \gets f_\theta(h_t + \tfrac{1}{2}k_2, t + \tfrac{1}{2}\Delta t)$
\State $k_4 \gets f_\theta(h_t + k_3, t + \Delta t)$
\State $h_{t+1} \gets h_t + \frac{\Delta t}{6}(k_1 + 2k_2 + 2k_3 + k_4)$ \Comment{combinación ponderada}
\end{algorithmic}

En lugar de las conexiones residuales simples de Euler $h_{t+1} = h_t + f_\theta(h_t)$, el bloque RK4 calcula cuatro evaluaciones intermedias y las combina con los pesos clásicos de RK4. Para evitar gradientes que se desvanecen, la arquitectura introduce \textbf{compuertas aprendidas} que interpolar entre aproximaciones intermedias:

\begin{algorithmic}[1]
\State $g \gets \sigma(\text{Lineal}([k_1, k_2, k_3, k_4]))$ \Comment{compuerta aprendible}
\State $h_{t+1} \gets h_t + g \cdot k_1 + (1-g) \cdot k_2$ \Comment{refinamiento interpolado}
\end{algorithmic}

Esta formulación reduce el número efectivo de parámetros mediante el uso compartido de pesos---la misma $f_\theta$ se evalúa múltiples veces---proporcionando un refinamiento de trayectoria más rico.

\begin{itemize}[leftmargin=*]
    \item \textbf{Complejidad de entrenamiento}: $\mathcal{O}(k \cdot n^2 \cdot d)$ donde $k$ es el orden RK (4 para RK4).
    \item \textbf{Complejidad de inferencia}: $\mathcal{O}(k \cdot n)$ por token (mayor sobrecarga constante).
    \item \textbf{Fortalezas}: Precisión significativamente mayor en tareas de generación (BLEU de última generación); eficiente en parámetros mediante uso compartido de pesos.
    \item \textbf{Debilidades}: Alto costo de cómputo por paso; la latencia de inferencia aumenta por un factor constante $k$; requiere compuertas complejas para la estabilidad.
\end{itemize}

\subsection{Memoria en Tiempo de Prueba: Titans}

La arquitectura Titans \citep{behrouz_titans_2025} introduce una dimensión ortogonal: en lugar de pesos estáticos, el modelo mantiene una memoria aprendible que se \textit{actualiza durante la inferencia} basada en una señal impulsada por sorpresa:

\begin{algorithmic}[1]
\State \textbf{Atención local a corto plazo:} Aplicar atención estándar o dispersa dentro de una ventana de contexto fija $c$
\State $y_t^{\text{local}} \gets \text{Atencion}(q_t, k_{[t-c:t]}, v_{[t-c:t]})$
\State \textbf{Señal de actualización de memoria (sorpresa):}
\State $S_t \gets \eta_t S_{t-1} - \theta_t \nabla_\ell(\mathcal{M}_{t-1}; x_t)$ \Comment{sorpresa como gradiente de la pérdida respecto a la memoria}
\State $\mathcal{M}_t \gets (1 - \alpha_t) \mathcal{M}_{t-1} + S_t$ \Comment{memoria actualizada mediante EMA}
\State \textbf{Recuperación de memoria a largo plazo:}
\State $y_t^{\text{memoria}} \gets \mathcal{M}_t^* (q_t)$ \Comment{consultar módulo de memoria para información recuperada}
\State \textbf{Combinación de salida:}
\State $y_t \gets y_t^{\text{local}} + \text{compuerta}(y_t^{\text{memoria}})$ \Comment{combinar memoria local y recuperada}
\end{algorithmic}

El módulo de memoria $\mathcal{M}$ se actualiza literalmente durante el paso hacia adelante calculando gradientes de una pérdida asociativa y aplicando pasos de SGD con momento. Las tasas de decaimiento $\eta_t, \theta_t, \alpha_t$ son en sí mismas dependientes de la entrada, permitiendo al modelo cambiar de paradigma de memoria cuando el contexto cambia.

\begin{itemize}[leftmargin=*]
    \item \textbf{Complejidad de entrenamiento}: Aproximadamente $\mathcal{O}(c^2 + n \cdot f)$ donde $c$ es la ventana local y $f$ es la sobrecarga de actualización de memoria.
    \item \textbf{Complejidad de inferencia}: $\mathcal{O}(c^2)$ atención local más $\mathcal{O}(1)$ recuperación de memoria por token.
    \item \textbf{Fortalezas}: Maneja longitudes de contexto extremas; permite recuperación asociativa real; la memoria se adapta a la entrada.
    \item \textbf{Debilidades}: Significativamente más complejo; la inferencia incluye cálculo de gradientes; mayor sobrecarga de coordinación.
\end{itemize}

\subsection{Atención Latente Multi-Cabeza (MLA)}
La Atención Latente Multi-Cabeza (MLA, por sus siglas en inglés) \citep{mehta_mla_2025} comprime el estado clave--valor por token en un único vector latente de rango bajo $c_{KV}$ de rango $r_{kv}$, el cual es la única cantidad materializada en el caché KV durante la decodificación. Un par de matrices de up-proyección reconstruye las claves y valores completos de todas las cabezas al vuelo, mientras que RoPE se aplica a los tensores de consulta y clave descomprimidos para preservar la estructura posicional. Los autores muestran que el ancho latente óptimo de Pareto se sitúa cerca de $r_{kv} = d_{\text{oculto}} / 2$, punto en el cual MLA reduce a la mitad el caché KV pero iguala la calidad de la atención multi-cabeza (MHA) en modelos pequeños.

\subsubsection{Formulación Matemática}
\[
c_{KV} = W_{DKV}\, x \in \mathbb{R}^{r_{kv}}, \qquad
k = W_{UK}\, c_{KV}, \qquad
v = W_{UV}\, c_{KV}, \qquad
q = W_{Q}\, x.
\]
\[
\tilde q = \text{RoPE}(q), \quad \tilde k = \text{RoPE}(k), \qquad
\text{atencion} = \texttt{softmax}\!\left(\frac{\tilde q\, \tilde k^\top}{\sqrt{d_h}}\right) v.
\]

\subsubsection{Pseudocódigo Algorítmico}
\begin{algorithmic}[1]
\State \textbf{Entrada:} estado oculto $x \in \mathbb{R}^{n \times d}$, rango $r_{kv}$, proyecciones $W_{DKV}, W_{UK}, W_{UV}, W_{Q}$
\State $c_{KV} \gets W_{DKV}\, x$ \Comment{down-proyección a latente de rango $r_{kv}$ (en caché en inferencia)}
\State $k \gets W_{UK}\, c_{KV}$, \quad $v \gets W_{UV}\, c_{KV}$ \Comment{up-proyección a cabezas completas}
\State $q \gets W_{Q}\, x$
\State $\tilde q \gets \text{RoPE}(q)$, \quad $\tilde k \gets \text{RoPE}(k)$ \Comment{aplicar embeddings rotatorios}
\State $\text{pesos} \gets \texttt{softmax}(\tilde q\, \tilde k^\top / \sqrt{d_h})$
\State \textbf{Salida:} $\mathbf{Y} = \text{pesos} \cdot v$
\end{algorithmic}

\subsubsection{Características Clave}
\begin{itemize}[leftmargin=*]
    \item \textbf{Complejidad de entrenamiento}: $\mathcal{O}(n^2 \cdot d)$ (atención completa sobre cabezas descomprimidas).
    \item \textbf{Complejidad de inferencia}: $\mathcal{O}(n)$ por token, espacio de caché KV $\mathcal{O}(r_{kv} \cdot n)$ (sólo latente).
    \item \textbf{Entrenable}: totalmente entrenable; $r_{kv}$ es un hiperparámetro (óptimo de Pareto $\approx d/2$).
    \item \textbf{Fortalezas}: reduce a la mitad el caché KV frente a MHA con calidad equivalente; compatible con RoPE; cuello de botella latente limpio.
    \item \textbf{Limitaciones}: la up-proyección añade FLOPs de decodificación; el rango latente es un presupuesto global fijo compartido entre cabezas.
\end{itemize}

\subsection{Atención Latente por Consultas Agrupadas (GQLA)}
La Atención Latente por Consultas Agrupadas (GQLA, por sus siglas en inglés) \citep{meng_gqla_2026} es una modificación mínima de MLA que expone dos rutas de decodificación algebraicamente equivalentes sobre los mismos pesos entrenados: una ruta MQA-absorb (que fusiona la up-proyección en la consulta, óptima en hardware tipo H100) y una ruta GQA-expanded (que materializa claves/valores completos, óptima en hardware H20 / predicción multi-token). Las dos rutas son matemáticamente idénticas porque la up-proyección es lineal, de modo que $\texttt{softmax}(q_{\text{absorbido}} \cdot c_{KV}^\top) = \texttt{softmax}((q \cdot W_{UK}^\top) \cdot c_{KV}^\top)$, permitiendo un despacho adaptativo al hardware sin reentrenamiento.

\subsubsection{Formulación Matemática}
Mismo latente que MLA:
\[
c_{KV} = W_{DKV}\, x, \qquad k = W_{UK}\, c_{KV}, \qquad v = W_{UV}\, c_{KV}, \qquad q = W_{Q}\, x.
\]
\[
\text{Ruta A (MQA-absorb): } \text{atencion} = \texttt{softmax}\!\left(\frac{(q\, W_{UK}^\top)\, c_{KV}^\top}{\sqrt{d_h}}\right) (W_{UV}\, c_{KV}).
\]
\[
\text{Ruta B (GQA-expanded): } \text{atencion} = \texttt{softmax}\!\left(\frac{q\, k^\top}{\sqrt{d_h}}\right) v, \quad k = W_{UK} c_{KV},\ v = W_{UV} c_{KV}.
\]

\subsubsection{Pseudocódigo Algorítmico}
\begin{algorithmic}[1]
\State \textbf{Entrada:} $x$, pesos $W_{DKV}, W_{UK}, W_{UV}, W_{Q}$, bandera de hardware $\text{hw} \in \{\text{H100}, \text{H20}\}$
\State $c_{KV} \gets W_{DKV}\, x$ \Comment{latente compartido, en caché}
\State $q \gets W_{Q}\, x$
\If{$\text{hw} == \text{H100}$} \Comment{ruta MQA-absorb}
    \State $q_{\text{abs}} \gets q\, W_{UK}^\top$ \Comment{absorber up-proyección en la consulta}
    \State $\text{pesos} \gets \texttt{softmax}(q_{\text{abs}}\, c_{KV}^\top / \sqrt{d_h})$
    \State $\mathbf{Y} \gets \text{pesos} \cdot (W_{UV}\, c_{KV})$
\Else \Comment{ruta GQA-expanded}
    \State $k \gets W_{UK}\, c_{KV}$, \quad $v \gets W_{UV}\, c_{KV}$
    \State $\text{pesos} \gets \texttt{softmax}(q\, k^\top / \sqrt{d_h})$
    \State $\mathbf{Y} \gets \text{pesos} \cdot v$
\EndIf
\State \textbf{Salida:} $\mathbf{Y}$
\end{algorithmic}

\subsubsection{Características Clave}
\begin{itemize}[leftmargin=*]
    \item \textbf{Complejidad de entrenamiento}: $\mathcal{O}(n^2 \cdot d)$, idéntica a MLA.
    \item \textbf{Complejidad de inferencia}: $\mathcal{O}(n)$ por token, estado latente $\mathcal{O}(r_{kv} \cdot n)$; ruta elegida por hardware.
    \item \textbf{Entrenable}: se entrena una sola vez; las dos rutas de decodificación reutilizan los mismos pesos.
    \item \textbf{Fortalezas}: decodificación adaptativa al hardware sin pérdida de calidad; la equivalencia algebraica garantiza consistencia.
    \item \textbf{Limitaciones}: mismas compensaciones de rango latente que MLA; la selección de ruta es una heurística de despliegue.
\end{itemize}

\subsection{Atención Multi-Cabeza de Rango Bajo (MLRA)}
La Atención Multi-Cabeza de Rango Bajo (MLRA, por sus siglas en inglés) \citep{liu_mlra_2026} refactoriza el latente único de MLA en $L$ sub-cabezas disjuntas cada una de rango $r/L$, de modo que el caché KV puede particionarse entre $L$ dispositivos para una decodificación tensor-paralela de 4 vías (cada dispositivo carga sólo $1/L$ del caché). Las up-proyecciones por sub-cabeza reconstruyen la porción correspondiente de las cabezas completas y se concatenan, logrando una aceleración de decodificación de 2.8$\times$ sobre MLA preservando el presupuesto latente global.

\subsubsection{Formulación Matemática}
\[
c_{KV} = [c_1, c_2, \ldots, c_L], \quad c_i \in \mathbb{R}^{r/L}, \qquad c_{KV} = W_{DKV}\, x.
\]
\[
k_i = W_{UK}^{(i)}\, c_i, \quad v_i = W_{UV}^{(i)}\, c_i, \qquad
k = [k_1; \ldots; k_L], \quad v = [v_1; \ldots; v_L].
\]
\[
\text{atencion} = \texttt{softmax}\!\left(\frac{q\, k^\top}{\sqrt{d_h}}\right) v.
\]

\subsubsection{Pseudocódigo Algorítmico}
\begin{algorithmic}[1]
\State \textbf{Entrada:} $x$, número de sub-cabezas $L$, proyecciones por sub-cabeza $W_{UK}^{(i)}, W_{UV}^{(i)}$
\State $c_{KV} \gets W_{DKV}\, x = [c_1, \ldots, c_L]$ \Comment{particionar latente en $L$ bloques}
\For{$i = 1, \ldots, L$ (en paralelo entre $L$ dispositivos)}
    \State $k_i \gets W_{UK}^{(i)}\, c_i$, \quad $v_i \gets W_{UV}^{(i)}\, c_i$
\EndFor
\State $k \gets \texttt{concat}(k_1, \ldots, k_L)$, \quad $v \gets \texttt{concat}(v_1, \ldots, v_L)$
\State $\text{pesos} \gets \texttt{softmax}(q\, k^\top / \sqrt{d_h})$
\State \textbf{Salida:} $\mathbf{Y} = \text{pesos} \cdot v$
\end{algorithmic}

\subsubsection{Características Clave}
\begin{itemize}[leftmargin=*]
    \item \textbf{Complejidad de entrenamiento}: $\mathcal{O}(n^2 \cdot d)$ (atención completa sobre cabezas concatenadas).
    \item \textbf{Complejidad de inferencia}: $\mathcal{O}(n)$ por token, estado $\mathcal{O}(r_{kv} \cdot n)$ particionado entre $L$ dispositivos.
    \item \textbf{Entrenable}: totalmente entrenable; $L$ controla la granularidad de partición.
    \item \textbf{Fortalezas}: aceleración de decodificación de 2.8$\times$; particionamiento limpio tensor-paralelo de 4 vías del caché; preserva el presupuesto latente total.
    \item \textbf{Limitaciones}: requiere que $L$ divida a $r_{kv}$; más matrices de proyección que MLA; all-reduce inter-dispositivo en la salida.
\end{itemize}

\subsection{Atención Tucker}
La Atención Tucker \citep{klein_tucker_2026} unifica MHA, GQA y MLA bajo una única factorización estilo Tucker de los tensores de pesos de consulta/clave/valor con rangos independientes $q_{\text{rank}}, k_{\text{rank}}, v_{\text{rank}}$. Se recupera MHA cuando todos los rangos igualan el tamaño oculto, GQA cuando $k_{\text{rank}} = v_{\text{rank}} < d$, y MLA cuando $k_{\text{rank}} = v_{\text{rank}}$ es igual al rango latente. Como la factorización es una repas lineal de pesos, el método es totalmente compatible con FlashAttention y RoPE, y produce un orden de magnitud menos parámetros para métricas comparables.

\subsubsection{Formulación Matemática}
\[
\mathbf{Q} = W_{Q,\text{core}}\, W_{Q,\text{factor}}\, x, \qquad
\mathbf{K} = W_{K,\text{core}}\, W_{K,\text{factor}}\, x, \qquad
\mathbf{V} = W_{V,\text{core}}\, W_{V,\text{factor}}\, x.
\]
\[
\text{atencion} = \texttt{softmax}\!\left(\frac{\mathbf{Q}\, \mathbf{K}^\top}{\sqrt{d_h}}\right) \mathbf{V}.
\]
Casos especiales: MHA $\Leftrightarrow q_{\text{rank}} = k_{\text{rank}} = v_{\text{rank}} = d$; GQA $\Leftrightarrow k_{\text{rank}} = v_{\text{rank}} < d$; MLA $\Leftrightarrow k_{\text{rank}} = v_{\text{rank}} = r_{kv}$.

\subsubsection{Pseudocódigo Algorítmico}
\begin{algorithmic}[1]
\State \textbf{Entrada:} $x$, rangos $q_{\text{rank}}, k_{\text{rank}}, v_{\text{rank}}$, matrices core/factor
\State $\mathbf{Q} \gets W_{Q,\text{core}}\, W_{Q,\text{factor}}\, x$ \Comment{consulta factorizada de rango $q_{\text{rank}}$}
\State $\mathbf{K} \gets W_{K,\text{core}}\, W_{K,\text{factor}}\, x$ \Comment{clave factorizada de rango $k_{\text{rank}}$}
\State $\mathbf{V} \gets W_{V,\text{core}}\, W_{V,\text{factor}}\, x$ \Comment{valor factorizado de rango $v_{\text{rank}}$}
\State \textbf{Aplicar RoPE a $\mathbf{Q}, \mathbf{K}$ si se usa codificación posicional}
\State $\text{pesos} \gets \texttt{softmax}(\mathbf{Q}\, \mathbf{K}^\top / \sqrt{d_h})$ \Comment{compatible con FlashAttention}
\State \textbf{Salida:} $\mathbf{Y} = \text{pesos} \cdot \mathbf{V}$
\end{algorithmic}

\subsubsection{Características Clave}
\begin{itemize}[leftmargin=*]
    \item \textbf{Complejidad de entrenamiento}: $\mathcal{O}(n^2 \cdot d)$ sobre los tensores reconstruidos.
    \item \textbf{Complejidad de inferencia}: $\mathcal{O}(n)$ por token; tamaño de caché gobernado por $k_{\text{rank}}$.
    \item \textbf{Entrenable}: los rangos son hiperparámetros; la factorización es una repas de pesos directa.
    \item \textbf{Fortalezas}: unifica MHA/GQA/MLA; reducción de parámetros de un orden de magnitud con métricas equivalentes; compatible con FlashAttention/RoPE.
    \item \textbf{Limitaciones}: dos matmuls por proyección añaden latencia; la selección de rango es un problema de ajuste por tarea.
\end{itemize}

\subsection{Atención de Cabezas Entrelazadas (IHA)}
La Atención de Cabezas Entrelazadas (IHA, por sus siglas en inglés) \citep{duvvuri_iha_2026} construye $P$ pseudo-cabezas por cabeza original (típicamente $P = H$) donde cada consulta/clave/valor pseudo es una combinación lineal aprendida de todas las $H$ cabezas originales. Esto induce hasta $P^2$ patrones de atención distintos por cabeza con una sobrecarga de $\mathcal{O}(H^2 P)$. Teóricamente, IHA necesita sólo $\Theta(\sqrt{k}\, n^2)$ parámetros frente a $\Theta(k\, n^2)$ de MHA en la tarea Polinomial; empíricamente entrega $+10$--$20\%$ en RULER de recuperación multi-clave, $+5.8\%$ en GSM8K y $+2.8\%$ en MATH-500.

\subsubsection{Formulación Matemática}
\[
q_{\text{pseudo}}[p] = \sum_{h=1}^{H} \text{mix}_q[p, h]\, q[h], \quad
k_{\text{pseudo}}[p] = \sum_{h} \text{mix}_k[p, h]\, k[h], \quad
v_{\text{pseudo}}[p] = \sum_{h} \text{mix}_v[p, h]\, v[h].
\]
\[
\text{atencion}_p = \texttt{softmax}\!\left(\frac{q_{\text{pseudo}}[p]\, k_{\text{pseudo}}[p]^\top}{\sqrt{d_h}}\right) v_{\text{pseudo}}[p],
\qquad
\mathbf{Y}[h] = \frac{1}{P} \sum_{p} \text{atencion}_p[h].
\]

\subsubsection{Pseudocódigo Algorítmico}
\begin{algorithmic}[1]
\State \textbf{Entrada:} por cabeza $q[h], k[h], v[h]$, matrices de mezcla $\text{mix}_q, \text{mix}_k, \text{mix}_v \in \mathbb{R}^{P \times H}$
\For{$p = 1, \ldots, P$}
    \State $q_p \gets \sum_h \text{mix}_q[p,h]\, q[h]$ \Comment{mezcla aprendida entre cabezas}
    \State $k_p \gets \sum_h \text{mix}_k[p,h]\, k[h]$, \quad $v_p \gets \sum_h \text{mix}_v[p,h]\, v[h]$
    \State $\text{atencion}_p \gets \texttt{softmax}(q_p\, k_p^\top / \sqrt{d_h})\, v_p$
\EndFor
\State \textbf{Salida:} $\mathbf{Y}[h] \gets \frac{1}{P} \sum_p \text{atencion}_p[h]$ \Comment{promediar salidas pseudo-cabeza por cabeza}
\end{algorithmic}

\subsubsection{Características Clave}
\begin{itemize}[leftmargin=*]
    \item \textbf{Complejidad de entrenamiento}: $\mathcal{O}(n^2 \cdot H \cdot P)$ con $P \approx H$, es decir $\mathcal{O}(n^2 \cdot H^2)$.
    \item \textbf{Complejidad de inferencia}: $\mathcal{O}(n)$ por token con caché KV; sobrecarga extra de mezcla $\mathcal{O}(H^2 P)$.
    \item \textbf{Entrenable}: las matrices de mezcla se aprenden; $P$ es un hiperparámetro.
    \item \textbf{Fortalezas}: eficiencia de parámetros $\Theta(\sqrt{k})$ en Polinomial; grandes ganancias en recuperación y razonamiento.
    \item \textbf{Limitaciones}: costo de mezcla cuadrático en cabezas; más patrones de atención a computar; latencia añadida con $H$ grande.
\end{itemize}

\subsection{Atención laTenT por Cabezas Agrupadas (GTA)}
La Atención laTenT por Cabezas Agrupadas (GTA, por sus siglas en inglés) \citep{sun_gta_2025} combina dos compresiones: (1) un \emph{mapa de atención compartido} --- un tensor de puntajes de atención por grupo de cabezas, reutilizado entre todas las cabezas del grupo, reduciendo el caché de claves; y (2) un \emph{decodificador de valores no lineal} --- el caché de valores se comprime a un latente mediante una down-proyección y se reconstruye por una up-proyección con compuerta SiLU antes de la proyección de salida. GTA reduce los FLOPs de atención hasta 62.5\% frente a GQA, encoge el caché KV hasta 70\%, y produce una aceleración de inferencia de 2$\times$ de extremo a extremo.

\subsubsection{Formulación Matemática}
\[
q_g = \texttt{mean}(q[\text{grupo } g]), \quad k_g = \texttt{mean}(k[\text{grupo } g]), \quad
\text{atencion}_g = \texttt{softmax}\!\left(\frac{q_g\, k_g^\top}{\sqrt{d_h}}\right) \text{ (reutilizado en el grupo)}.
\]
\[
v_{\text{latente}} = W_{DV}\, x, \qquad v = \text{SiLU}(W_{UV}\, v_{\text{latente}}), \qquad
\mathbf{Y} = \text{atencion}_g \cdot v.
\]

\subsubsection{Pseudocódigo Algorítmico}
\begin{algorithmic}[1]
\State \textbf{Entrada:} $q, k$ por cabeza, oculto $x$, partición de grupos $\{g\}$
\State $v_{\text{latente}} \gets W_{DV}\, x$ \Comment{comprimir valores a latente (en caché)}
\State $v \gets \text{SiLU}(W_{UV}\, v_{\text{latente}})$ \Comment{decodificación no lineal de valores}
\For{cada grupo $g$}
    \State $q_g \gets \texttt{mean}(q[g])$, \quad $k_g \gets \texttt{mean}(k[g])$ \Comment{consultas/claves compartidas por grupo}
    \State $\text{atencion}_g \gets \texttt{softmax}(q_g\, k_g^\top / \sqrt{d_h})$
    \State $\mathbf{Y}[g] \gets \text{atencion}_g \cdot v[g]$ \Comment{un mapa reutilizado para todas las cabezas en $g$}
\EndFor
\State \textbf{Salida:} $\mathbf{Y}$
\end{algorithmic}

\subsubsection{Características Clave}
\begin{itemize}[leftmargin=*]
    \item \textbf{Complejidad de entrenamiento}: $\mathcal{O}(n^2 \cdot G \cdot d_h)$ donde $G$ es el número de grupos (frente a $H$ cabezas en GQA).
    \item \textbf{Complejidad de inferencia}: $\mathcal{O}(n)$ por token; caché de valores $\mathcal{O}(r_v \cdot n)$, caché de claves reducido por compartición de grupo.
    \item \textbf{Entrenable}: la asignación de grupo y el rango latente son entrenables/hiperparámetros.
    \item \textbf{Fortalezas}: hasta 62.5\% de reducción de FLOPs de atención, 70\% de encogimiento de caché KV, 2$\times$ de aceleración de extremo a extremo.
    \item \textbf{Limitaciones}: el mapa compartido reduce la diversidad por cabeza; el decodificador no lineal añade latencia; el agrupamiento debe elegirse con cuidado.
\end{itemize}

\subsection{Atención Latente Temporal Multi-Cabeza (MTLA)}
La Atención Latente Temporal Multi-Cabeza (MTLA, por sus siglas en inglés) \citep{deng_mtla_2025} extiende MLA a lo largo del eje temporal: una hyper-red fusiona dinámicamente $m$ vectores temporalmente adyacentes del caché KV en un único slot, reduciendo la longitud temporal por un factor de $m$. Una máscara causal consciente del stride mantiene el entrenamiento paralelo consistente con la inferencia autorregresiva, produciendo una aceleración de decodificación de 5.3$\times$ y una reducción de memoria GPU de 8.3$\times$ frente a MHA en traducción de voz En--De.

\subsubsection{Formulación Matemática}
\[
c_{KV}^{(t)} = W_{DKV}\, x^{(t)} \in \mathbb{R}^{r_{kv}} \text{ (latente por token)}.
\]
\[
\bar c_{KV}^{(j)} = \sum_{i=0}^{m-1} \alpha_i\, c_{KV}^{(s j + i)}, \qquad \alpha_i = \texttt{sigmoid}(\text{HyperNet}(c_{KV}^{(s j + i)})),
\]
donde $s$ es el stride y $m$ la ventana de fusión. Las consultas atienden sobre slots fusionados $\bar c_{KV}$ con una máscara causal consciente del stride $M_{\text{stride}}$:
\[
\text{atencion} = \texttt{softmax}\!\left(\frac{q\, \bar k^\top}{\sqrt{d_h}} + M_{\text{stride}}\right) \bar v, \quad \bar k = W_{UK}\, \bar c_{KV},\ \bar v = W_{UV}\, \bar c_{KV}.
\]

\subsubsection{Pseudocódigo Algorítmico}
\begin{algorithmic}[1]
\State \textbf{Entrada:} latentes por token $c_{KV}^{(t)}$, ventana de fusión $m$, stride $s$
\State \textbf{Fase de fusión (hyper-red):}
\For{cada slot fusionado $j = 0, 1, \ldots$}
    \State $\alpha_i \gets \texttt{sigmoid}(\text{HyperNet}(c_{KV}^{(s j + i)}))$ para $i = 0, \ldots, m-1$
    \State $\bar c_{KV}^{(j)} \gets \sum_{i=0}^{m-1} \alpha_i\, c_{KV}^{(s j + i)}$ \Comment{promedio con compuerta sobre la ventana}
\EndFor
\State $\bar k \gets W_{UK}\, \bar c_{KV}$, \quad $\bar v \gets W_{UV}\, \bar c_{KV}$
\State $\text{pesos} \gets \texttt{softmax}(q\, \bar k^\top / \sqrt{d_h} + M_{\text{stride}})$ \Comment{máscara causal consciente del stride}
\State \textbf{Salida:} $\mathbf{Y} = \text{pesos} \cdot \bar v$
\end{algorithmic}

\subsubsection{Características Clave}
\begin{itemize}[leftmargin=*]
    \item \textbf{Complejidad de entrenamiento}: $\mathcal{O}(n^2 \cdot d / m)$ sobre la secuencia fusionada (longitud $\lceil n/m \rceil$) más costo de la hyper-red.
    \item \textbf{Complejidad de inferencia}: $\mathcal{O}(n/m)$ por token sobre slots fusionados; estado $\mathcal{O}(r_{kv} \cdot n / m)$.
    \item \textbf{Entrenable}: la hyper-red y la ventana de fusión $m$ se aprenden/eliges; la máscara consciente del stride garantiza consistencia.
    \item \textbf{Fortalezas}: aceleración de decodificación de 5.3$\times$, reducción de memoria GPU de 8.3$\times$ frente a MHA; compresión temporal ortogonal a la compresión de rango.
    \item \textbf{Limitaciones}: la ventana de fusión puede difuminar detalle temporal fino; la hyper-red añade cómputo; el diseño de la máscara no es trivial.
\end{itemize}

\subsection{Comparación y Síntesis Arquitectónica}

\small
\begin{longtable}{p{2.0cm}p{1.8cm}p{1.6cm}p{1.6cm}p{4.5cm}}
\toprule
\textbf{Arquitectura} & \textbf{Entrenar} & \textbf{Inferir} & \textbf{Estado} & \textbf{Característica Clave} \\
\midrule
\endhead
Atención Estándar & $O(n^2 d)$ & $O(n)$ caché & $O(nd)$ & Expresividad perfecta, enrutamiento completo de tokens, cuello de botella KV. \\
GQA & $O(n^2 d)$ & $O(n)$ caché & $O(h_k \cdot n \cdot d_h)$ & Compresión KV ajustable, probado a escala, reducción de caché en factor $G \times$. \\
Atención Sigmoide & $O(n^2 d)$ & $O(n)$ caché & $O(nd)$ & Compuerta elemento a elemento, eficiente en hardware, estable a escala con norma híbrida. \\
RetNet & $O(n^2 d)$ & $O(1)$ & $O(d^2)$ & Decaimiento multi-escala, triple modo de cómputo, patrón de olvido fijo. \\
Mamba & $O(nd)$ & $O(1)$ & $O(d)$ & Dinámica selectiva dependiente de la entrada, entrenamiento e inferencia lineales, escaneo de hardware. \\
Transformer EDO & $O(kn^2 d)$ & $O(kn)$ & $O(nd)$ & Refinamiento por integración numérica, uso compartido de pesos por etapas, mayor precisión. \\
Titans & $O(c^2 + nf)$ & $O(c^2 + 1)$ & $O(1)$ & Adaptación de memoria en tiempo de prueba, contexto extremo, actualizaciones de parámetros durante inferencia. \\
MLA & $O(n^2 d)$ & $O(n)$ caché & $O(r_{kv} n)$ & KV latente de rango $r_{kv}\!\approx\!d/2$, reduce a la mitad el caché, iguala calidad MHA. \\
GQLA & $O(n^2 d)$ & $O(n)$ caché & $O(r_{kv} n)$ & Mismos pesos MLA, dos rutas de decodificación algebraicamente equivalentes adaptativas al hardware. \\
MLRA & $O(n^2 d)$ & $O(n)$ caché & $O(r_{kv} n)$ particionado & Latente dividido en $L$ sub-cabezas para decodificación tensor-paralela de 4 vías, 2.8$\times$ de aceleración. \\
Tucker & $O(n^2 d)$ & $O(n)$ caché & $O(k_{\text{rank}} n)$ & Factorización Tucker unifica MHA/GQA/MLA, un orden de magnitud menos parámetros. \\
IHA & $O(n^2 H^2)$ & $O(n)$ caché & $O(H n d_h)$ & Pseudo-cabezas como mezclas aprendidas de cabezas, $\Theta(\sqrt{k})$ parámetros, +10--20\% RULER. \\
GTA & $O(n^2 G d_h)$ & $O(n)$ caché & $O(r_v n)$ & Mapa compartido por grupo + decodificador de valores SiLU, 62.5\% menos FLOPs, 70\% menos caché. \\
MTLA & $O(n^2 d / m)$ & $O(n/m)$ & $O(r_{kv} n/m)$ & Fusión temporal de $m$ slots KV, 5.3$\times$ decodificación, 8.3$\times$ reducción de memoria. \\
\bottomrule
\end{longtable}
\normalsize

El campo exhibe una progresión clara a lo largo de dos ejes: (i) \textbf{eficiencia computacional}, pasando de $O(n^2)$ a $O(n)$ u $O(1)$, y (ii) \textbf{adaptabilidad de memoria}, pasando de pesos estáticos a modelos dinámicos actualizados en tiempo de prueba. La atención estándar sigue siendo la línea base de expresividad; Mamba y RetNet representan la frontera práctica de eficiencia; Titans introduce un eje de innovación ortogonal (aprendizaje en tiempo de prueba). La elección de arquitectura refleja la restricción fundamental de ingeniería: expresividad versus costo de despliegue.

